\chapter{显存占用严谨核算与 GPU 算力 Roofline 实测}

\section{Lab 08：显存开销与 Roofline 性能模型}

在进行大模型推理系统设计或显卡选型时，有两大黄金法则：
\begin{enumerate}
  \item \textbf{显存容量法则}：决定了你「能不能跑得下」（模型参数 + KV Cache + 运行时开销 $\le$ 显存容量）。
  \item \textbf{Roofline 性能模型与访存带宽法则}：决定了你「能跑多快」（自回归解码受限于 GPU 显存带宽）。
\end{enumerate}

\vspace{0.8em}\hrule\vspace{0.8em}

\section{[推导] 1. 显存开销组成详解}

大模型在推理运行时的显存占用主要由以下四部分组成：

\[
\text{Total Memory} = \text{Mem}_{\text{weights}} + \text{Mem}_{\text{KV\_Cache}} + \text{Mem}_{\text{activation}} + \text{Mem}_{\text{CUDA\_overhead}}
\]

\subsection{(1) 模型权重 (Model Weights)}

\begin{itemize}
  \item \textbf{FP16 / BF16}: 每个参数占用 \textbf{2 字节 (Bytes)}。
  \item 例如 7B 模型需要约 $7 \times 2 = 14 \text{ GB}$。
  \item \textbf{INT8 / FP8}: 每个参数占用 \textbf{1 字节}（7B 约 7 GB）。
  \item \textbf{INT4 / AWQ / GPTQ}: 每个参数占用 \textbf{0.5 字节}（7B 约 3.5 GB）。
\end{itemize}

\subsection{(2) KV Cache 显存}

这是动态显存开销的最大头，随并发请求数和序列长度成正比膨胀：
对于采用分组查询注意力（GQA, Grouped-Query Attention）的模型：
\begin{itemize}
  \item 层数 $L$
  \item KV 头数 $n_{\text{kv\_heads}}$，单头维度 $d_{\text{head}}$
  \item 并发请求数 $B$（Batch Size）
  \item 序列总长度 $S$（Prompt 长度 + 已生成长度）
  \item 存储精度字节数 $b$（FP16 为 2）
\end{itemize}

\[
\text{KV Cache (Bytes)} = 2 \times L \times (n_{\text{kv\_heads}} \times d_{\text{head}}) \times S \times B \times b
\]

\begin{tipBox}[核心直觉与思考]
{}[思考] \textbf{MHA vs GQA 的显存革命}：
- 早期模型（如 LLaMA-1 65B）使用 Multi-Head Attention (MHA)，KV 头数等于 Query 头数（例如 64 头）。
- 现代模型（如 Llama-3、Qwen2.5）几乎全部采用 GQA（例如 8 个 KV 头），KV Cache 显存直接暴降为原来的 $1/4 \sim 1/8$！
\end{tipBox}

\subsection{(3) 激活值 (Activations)}

在 Prefill 阶段较大（因为输入序列长），但在 Decode 阶段因为一次只输入 1 个 token，激活值极小（通常仅几十 MB 到几百 MB）。

\subsection{(4) CUDA 上下文与运行时开销}

PyTorch / CUDA runtime、NCCL 通信缓冲区等，通常占用约 $1 \sim 2 \text{ GB}$。

\vspace{0.8em}\hrule\vspace{0.8em}

\section{[高速] 2. Roofline 模型与为什么 Decode 是访存受限？}

\subsection{计算密度（Arithmetic Intensity）}

\[
\text{Arithmetic Intensity (AI)} = \frac{\text{计算量 (FLOPs)}}{\text{访存数据量 (Bytes)}}
\]

GPU 的极限速度受制于：
\begin{itemize}
  \item \textbf{最大算力 (TFLOPs)}
  \item \textbf{显存带宽 (GB/s)}
\end{itemize}

两者的交界点为硬件的转折算术密度 $AI_{\text{roof}} = \frac{\text{Peak FLOPs}}{\text{Peak Bandwidth}}$。

\begin{center}
\small
\begin{tabularx}{\textwidth}{l X X X X X }
\toprule
\textbf{阶段} & \textbf{输入规模} & \textbf{计算量} & \textbf{访存量} & \textbf{算术强度 $AI$} & \textbf{性能瓶颈} \\
\midrule
\textbf{Prefill 阶段} & $N$ 个 token 并行 & $O(N \cdot P)$ 大矩阵乘法 & 读取一次模型权重 $P$ & 很高 ($> 100$) & \textbf{算力受限 (Compute-Bound)}，GPU 计算核心满载 \\
\textbf{Decode 阶段} & 每次仅 1 个 token & 仅对 1 个 token 做矩阵-向量乘 & 必须把全部模型权重和全部历史 KV Cache 从显存重新读一遍！ & 极低 ($\approx 1 \sim 2$) & \textbf{访存受限 (Memory-Bound)}，算力被闲置，卡在显存带宽搬运上 \\
\bottomrule
\end{tabularx}
\end{center}

\subsection{单 Token 解码耗时下界理论估算}

在 Batch Size = 1 时，生成 1 个 token 必须把整个模型权重读取一遍：

\[
T_{\text{decode\_min}} \approx \frac{\text{模型权重体积 (Bytes)}}{\text{显存有效读取带宽 (Bytes/s)}}
\]

\begin{noteBox}[核心导读与背景]
\textbf{例如}：在显存带宽为 448 GB/s 的显卡上运行 14 GB 的 7B FP16 模型：
每生成 1 个 token 至少需要 $14 \text{ GB} / 448 \text{ GB/s} \approx 31.25 \text{ ms}$（即吞吐上限约为 $32 \text{ tokens/s}$）。
如果通过 INT4 量化将权重压缩为 3.5 GB，理论下界降为 $3.5 / 448 \approx 7.8 \text{ ms}$（吞吐暴涨到 $128 \text{ tokens/s}$）！这就是量化在推理阶段能带来巨大速度提升的根本原因！

{}[注意] 上面的 448 GB/s 是\textbf{厂商标称值}，实际达不到（本机实测约 380 GB/s，见下面"一次实测结果"）。
用实测值算，同一个例子的下界是 $14/380 \approx 36.8 \text{ ms}$ 而不是 31.25 ms —— \textbf{标称值会乐观约 18\%}。
\texttt{memory\_calculator.py} 现在两种都打印，默认用实测值。
\end{noteBox}

\begin{noteBox}[核心导读与背景]
\textbf{单位口径}：本节算例用\textbf{十进制 GB}（7B $\times$ 2 字节 = $14 \times 10^9$ B = 14 GB），
因为要和 GB/s 相除。而 \texttt{memory\_calculator.py} 的权重/KV 打印用 \textbf{GiB}（除以 $1024^3$），
因为要和显卡容量比较。同一个 7B FP16 模型：14 GB = 13.04 GiB，\textbf{两个都对，差 7.4\%}。
混用是这一行最常见的错误，脚本里专门有一段注释说明。
\end{noteBox}

\vspace{0.8em}\hrule\vspace{0.8em}

\section{[工具] 显存与性能计算工具使用}

运行 \textbf{\texttt{memory\_calculator.py}}：

\begin{lstlisting}[language=bash]
python3 lab08_memory_and_roofline/memory_calculator.py
\end{lstlisting}

可评估主流模型（Llama-3-8B/70B、Qwen-2.5-7B/72B 等）在各类 GPU（如 RTX 5060 Ti、RTX 4090、A100）下的显存与并发上限。

\vspace{0.8em}\hrule\vspace{0.8em}

\section{[实验] 在你自己的显卡上实测 Roofline（需要 PyTorch）}

理论推导见 \textbf{第 8 篇：硬件视角}。运行：

\begin{lstlisting}[language=bash]
.venv/bin/python lab08_memory_and_roofline/measure_gpu_roofline.py
\end{lstlisting}

它会测出你这张卡的真实带宽 $\beta$、BF16 峰值 $\pi$、拐点 AI* = $\pi$/$\beta$，并验证两个关键推导：
\begin{enumerate}
  \item \textbf{Decode 线性层的算术强度 $\approx$ batch size}：batch 从 1 到 16，耗时几乎不变；
  \item \textbf{Decode 注意力的算术强度 $\approx$ GQA 分组大小}：与 batch 无关，所以 GQA/MLA 不仅省显存，还直接提速。
\end{enumerate}

本机（RTX 5060 Ti 16GB）的一次实测结果：

\begin{noteBox}[核心导读与背景]
{}[计时] 以下计时类数字是本机某一次运行的\textbf{量级示例}，前提是显卡空闲、已预热；GPU 降频或有其它负载时可能慢 40\% 以上。请看\textbf{趋势和比例}，以你自己运行的输出为准。 例如有评审在并发负载下测到过 217\textasciitilde{}234 GB/s 的带宽。
\end{noteBox}

\begin{center}
\small
\begin{tabularx}{\textwidth}{p{3cm} p{4.5cm} X}
\toprule
\textbf{指标} & \textbf{厂商标称} & \textbf{实测} \\
\midrule
显存带宽 & 448 GB/s & $\approx$ 380 GB/s（拷贝），$\approx$ 400 GB/s（GEMV 读权重） \\
BF16 Tensor 稠密算力 & 宣传口径更高 & \textbf{$\approx$ 48 TFLOP/s} \\
拐点 AI* & — & \textbf{$\approx$ 130 FLOP/Byte}：batch < 100 左右的 Decode 都是访存受限 \\
\bottomrule
\end{tabularx}
\end{center}

\begin{warningBox}[避坑指南与核心警示]
{}[注意] 实测中还会看到 batch = 32/64 的线性层比 batch = 128 还慢——这是 cuBLAS 对这些形状选中了低效 kernel，是真实存在的现象，不是噪声。它提醒你：\textbf{Roofline 是上界，实现质量决定你离上界有多远。}
\end{warningBox}



\section{实验核心源码精选与解析}

本实验配套有完整可运行的工业级验证代码，重点实现细节如下：


\subsection{源码实现：\texttt{memory\_calculator.py}}

\lstinputlisting[language=Python, caption={\texttt{memory\_calculator.py}}]{code/lab08_memory_and_roofline_memory_calculator.py}


\subsection{源码实现：\texttt{measure\_gpu\_roofline.py}}

\lstinputlisting[language=Python, caption={\texttt{measure\_gpu\_roofline.py}}]{code/lab08_memory_and_roofline_measure_gpu_roofline.py}

